\chapter{Estado del arte}
\label{ch:estado-arte}

Este capítulo revisa las aproximaciones existentes al problema de construir y
mantener aplicaciones cuyo modelo de datos evoluciona, con especial atención a
la generación automática de interfaces y a su acoplamiento con el \gls{backend}.
El objetivo es situar la propuesta de este trabajo: identificar qué soluciones
existen en la actualidad, en qué medida no satisfacen por completo el objetivo
planteado y qué aporta el \gls{framework} desarrollado frente a ellas.

\section{Desarrollo rápido de aplicaciones web}
\label{sec:rad}

El concepto de \gls{rad} fue formalizado por James Martin en los años noventa~\cite{martin1991rad} como
una respuesta a los ciclos de desarrollo en cascada, demasiado rígidos para adaptarse a los
cambios de requisitos. En el contexto de las aplicaciones web modernas, el \gls{rad} se
manifiesta principalmente a través de tres estrategias:

\begin{enumerate}
  \item \textbf{Scaffolding:} generación automática de código \emph{boilerplate} a partir de un
    modelo de datos (p.~ej., Django Admin, Rails Scaffolding).
  \item \textbf{Frameworks declarativos:} el desarrollador describe \textit{qué} quiere
    (esquemas, configuraciones) y el \textit{framework} genera \textit{cómo} lograrlo.
  \item \textbf{Low-code / No-code:} plataformas con interfaces visuales que eliminan
    casi por completo la escritura de código (p.~ej., Retool, Appsmith).
\end{enumerate}

El enfoque adoptado en este proyecto se sitúa entre las estrategias 1 y 2: se parte de
modelos de datos explícitos (SQLAlchemy) y se genera automáticamente tanto el
\gls{api} REST como la interfaz de usuario, mediante una tubería de transformación de
esquemas.

\subsection{Comparativa con soluciones existentes}
\label{subsec:comparativa}

\begin{table}[H]
  \centering
  \caption{Comparativa de soluciones RAD web de código abierto.}
  \label{tab:comparativa-rad}
  \begin{tabular}{lcccc}
    \toprule
    \textbf{Solución} & \textbf{Backend} & \textbf{Frontend} & \textbf{Extensible} &
    \textbf{Schema-driven} \\
    \midrule
    Django Admin & Python & Integrado & Limitado & No \\
    Rails Scaffold & Ruby & ERB & Sí & No \\
    Retool & SaaS & Visual & Sí & Parcial \\
    Appsmith & Node.js & React & Sí & Parcial \\
    \textbf{uniback + ngt-gui} & \textbf{Python/FastAPI} & \textbf{Angular} &
    \textbf{Sí} & \textbf{Sí} \\
    \bottomrule
  \end{tabular}
\end{table}

La principal diferencia respecto a las soluciones existentes es que \textit{uniback} +
\textit{ngt-gui} están diseñados para integrarse en proyectos con arquitecturas propias,
sin imponer un \emph{stack} cerrado ni requerir una plataforma SaaS.

\section{Generación de interfaces dirigida por esquema}
\label{sec:soa-interfaces}

La idea de derivar interfaces de un esquema no es nueva. En el ecosistema web
existen \texttt{react-jsonschema-form} y \texttt{JSON Forms}, que renderizan
formularios a partir de \gls{jsonschema}. Estas soluciones, sin embargo, se
centran en el lado cliente y asumen que el esquema se proporciona ya
elaborado. En el lado servidor, herramientas como FastAPI exponen OpenAPI,
pero ese contrato describe \emph{operaciones}, no incluye pistas de
presentación, y no resuelve el acoplamiento del cliente al dominio.

La tabla~\ref{tab:soa} compara los enfoques. La aportación diferencial de este
trabajo no es inventar un renderizador de esquemas, sino \textbf{completar la cadena}:
derivar el esquema \emph{desde el ORM}, enriquecerlo con
pistas de presentación en el propio esquema, transportarlo por un contrato
uniforme, y consumirlo con una capa genérica que mantiene la compatibilidad
con el código heredado. Considero, además, de vital importancia la retrocompatibilidad, ya que
sería inútil proporcionar un software que requiera más coste de implementación que el ya existente.

\begin{table}[!htbp]
\caption{Comparación de enfoques de generación de interfaces.}
\label{tab:soa}
\centering
\begin{tabular}{|p{3.4cm}|c|c|c|c|}
\hline
\textbf{Enfoque} & \textbf{Desde ORM} & \textbf{Hints UI} &
\textbf{Contrato} & \textbf{Versionado} \\ \hline
react-jsonschema-form & No & Parcial & No & No \\ \hline
JSON Forms & No & Sí (uischema) & No & No \\ \hline
OpenAPI/Swagger UI & No & No & Operaciones & No \\ \hline
Admin autogenerados & Sí & Limitado & Acoplado & No \\ \hline
\textbf{Este trabajo} & \textbf{Sí} & \textbf{Sí (x-*)} &
\textbf{Uniforme} & \textbf{ETag} \\ \hline
\end{tabular}
\end{table}

Frente a los \emph{admin} autogenerados, que sí
parten del ORM pero acoplan la interfaz al servidor y son difíciles de
reutilizar fuera de su \emph{stack}, la solución propuesta mantiene frontend y
backend desacoplados mediante un contrato explícito y portable, que es
precisamente lo que habilita la reutilización en proyectos con modelos de
datos en evolución y con diferentes tecnologías de bases de datos.

\section{Patrones arquitectónicos relacionados}
\label{sec:patrones}

\subsection{Schema-first development}
\label{subsec:schema-first}

El enfoque \textit{schema-first} propone definir el contrato de datos antes de implementar
las capas de negocio y presentación~\cite{schemafirst}. En este proyecto el esquema no se
define manualmente sino que se \textbf{genera automáticamente} desde los modelos ORM,
combinando las ventajas del \textit{schema-first} con las del \textit{model-first}.

\subsection{Plugin pattern}
\label{subsec:plugin-pattern}

El sistema de plugins por capas implementado en \textit{uniback} sigue el patrón
\textbf{Chain of Responsibility}~\cite{gof}: cada plugin recibe el esquema producido por
la capa anterior, lo enriquece y lo pasa a la siguiente. Esto permite añadir o quitar
plugins sin modificar el núcleo del generador.

\subsection{CQRS y operaciones en bloque}
\label{subsec:cqrs}

El endpoint \texttt{/bulk} introduce una separación implícita entre comandos de escritura
(crear, actualizar, eliminar) y consultas de lectura, en la línea del patrón
\gls{crud}~\cite{fowler2002} extendido. El modo \textit{partial} (con savepoints SQL) y
el modo \textit{abort} (transacción global) ofrecen semántica transaccional configurable
por el cliente.
